\documentclass{article}
\input{~/studies/preamble}

\title{Diffusion Model}
\author{Hyoung Joon, Kim}
\date{\today}

\begin{document}
\maketitle
\tableofcontents
\newpage

\section{Introduction}
    \textbf{Diffusion models} or \emph{denoising diffusion probabilistic models}, recovers data from pure noise. The encoder gradually adds Gaussian noise to the given data,
    and make it into pure Gaussian noise. Then, the decoder learns how to recover the image from pure Gaussian noise.
\section{Main Idea}
    \subsection{Forward Encoder}
        \subsubsection{Diffusion Process}
            The \textbf{diffusion process} (forward process) adds Gaussian noise to the data gradually.
            \begin{defn}{Diffusion Process}{diffProcess}
                Let $\{\tb z_1,\dots,\tb z_T\}$ be the vectors of same size of given data $\tb x = \tb z_0$,
                and the process $\tb z_1,\dots,\tb z_T$ is defined as:
                \[
                \tb z_t = \sqrt{1-\beta_t}\tb z_{t-1} + \sqrt{\beta_t}\bm\epsilon_t
                \]
                where $0 < \beta < 1$ is the variance of the noise distribution and $\bm\epsilon_t \sim \mathcal{N}(\tb 0, \tb I)$. Notice that we can rewrite the equation in the following form;
                \[
                q(\tb z_t|\tb z_{t-1})=\mathcal{N}(\tb z_t|\sqrt{1-\beta_t}\tb z_{t-1}, \beta_t\tb I).
                \]
            \end{defn}
            These sequence of distributions forms Markov chain. The values of $\beta_1,\beta_2,\dots,\beta_T \in (0,1)$ are typically chosen by hand such that $\beta_1 < \beta_2 < \dots < \beta_T$.
        \subsubsection{Diffusion Kernel}
            Since the distributions $q(\tb z_t|\tb z_{t-1})$ are normal, the marginalized distribution $q(\tb z_t | \tb x)$ can be written in closed form:
            \begin{thm}{Diffusion Kernel}{diffKern}
                \[
                q(\tb z_t|\tb x) = \mathcal{N}(\tb z_t|\sqrt{\alpha_t}\tb x, (1-\alpha_t)\tb I)
                \]
                where $\alpha_t = \prod_{\tau=1}^{t}(1-\beta_\tau)$. This marginalized distribution is called \tb{diffusion kernel}. As the noises are added, in the limit $T \to \infty$,
                we have $q(\tb z_T|\tb x)=\mathcal{N}(\tb z_T|0,\tb I)$, and since the right-hand side is independent of $\tb x$, $q(\tb z_T|\tb x)=q(\tb z_T)$.
            \end{thm}
    \subsection{Reverse Decoder}
        \subsubsection{Conditional Distribution}
            Since the decoder learns how to denoise the corrupted data, it is natural to consider the probability distribution $q(\tb z_{t-1}|\tb z_t)=\frac{q(\tb z_t|\tb z_{t-1})q(\tb z_{t-1})}{q(\tb z_t)}$.
            The marginal distribution $q(\tb z_t)$ can be written as $q(\tb z_t)=\int q(\tb z_t|\tb x)p(\tb x)d\tb x$. Notice that it is intractable to evaluate the marginal distribution. So, we consider the conditional
            version of the reverse distribution conditioned on $\tb x$, which is given as:
            \begin{thm}{Reverse Conditional Distribution}{revCondDist}
                The conditional reversed distribution is given as:
                \begin{align}
                q(\tb z_{t-1}|\tb z_t, \tb x)&=\frac{q(\tb z_t|\tb z_{t-1},\tb x)q(\tb z_{t-1}|\tb x)}{q(\tb z_t|\tb x)}
                \\ &=\frac{q(\tb z_t|\tb z_{t-1})q(\tb z_{t-1}|\tb x)}{q(\tb z_t|\tb x)}
                \\ &=\mathcal{N}(\tb z_{t-1}|\tb m_t(\tb x, \tb z_t), \sigma_t^2\tb I)
                \end{align}
                where
                \begin{align*}
                    \tb m_t(\tb x, \tb z_t)&=\frac{(1-\alpha_{t-1})\sqrt{1-\beta_t}\tb z_t + \sqrt{\alpha_{t-1}}\beta_t\tb x}{1-\alpha_t}
                \\  \sigma_t^2&=\frac{\beta_t(1-\alpha_{t-1})}{1-\alpha_{t-1}}.
                \end{align*}
            \end{thm}
                We used Markov chain property $q(\tb z_t|\tb z_{t-1},\dots,\tb x) = q(\tb z_t|\tb z_{t-1})$ in (1) to (2), and used properties of Gaussian distributions in (2) to (3).
        \subsubsection{Reverse Process}
            The reverse process is modeled as Gaussian distribution:
            \begin{defn}{Reverse Process}{revProcess}
                Let $\beta_1,...\beta_T$ be given as the same schedule in Definition \ref{defn:diffProcess}. Then the reverse process is defined as
                \[
                p(\tb z_{t-1}|\tb z_t,\tb w)=\mathcal{N}(\tb z_{t-1}|\bm\mu(\tb z_t, \tb w, t),\beta_t\tb I)
                \]
                where $\bm\mu(\tb z_t,\tb w, t)$ is a deep neural network with parameters $\tb w$.
            \end{defn}
            Then the likelihood is given as:
            \begin{align*}
                p(\tb x| \tb w) &= \int p(\tb x,\tb z_1,\dots,\tb z_T|\tb w)d\tb z_1 \dots d\tb z_T
            \\  &= \int p(\tb z_T)\left\{ \prod_{t=2}^{T}p(\tb z_{t-1}|\tb z_t,\tb w) \right\}p(\tb x|\tb z_1,\tb w)d\tb z_1 \dots d\tb z_T
            \end{align*}
            and the objective is to maximize this likelihood. Here $p(\tb z_T)$ is assumed to equal $q(\tb z_T)$ and hence is given by $\mathcal{N}(\tb z_T|\bm 0,\tb I)$.
    \subsection{Training}
        \subsubsection{Evidence Lower Bound}
            Since the likelihood itself is intractable to be evaluated, we use ELBO like VAE to maximize the likelihood;
            \begin{align*}
                \log{p(\tb x|\tb w)}
                &=\mathcal{L}(\tb w) + \text{KL}(q(\tb z)\|p(\tb z|\tb x,\tb w))
            \\  \mathcal{L}
                &=\int q(\tb z)\log{\left\{ \frac{p(\tb x, \tb z|\tb w)}{q(\tb z)} \right\}}d\tb z
            \\  \text{KL}(q(\tb z)\|p(\tb z|\tb x,\tb w))
                &=-\int q(\tb z)\log{\left\{ \frac{p(\tb z|\tb x,\tb w)}{q(\tb z)} \right\}}d\tb z
            \end{align*}
            We now substitute $q(\tb z)$ and $p(\tb x,\tb z_1,\dots,\tb z_T|\tb w)$
            with
            \begin{align*}
                q(\tb z_1,\dots,\tb z_T|\tb x) &= q(\tb z_1|\tb x)\prod_{t=2}^{T}q(\tb z_t|\tb z_{t-1})
            \\  p(\tb x,\tb z_1,\dots,\tb z_T|\tb w) &= p(\tb z_t)\prod_{t=2}^{T}p(\tb z_{t-1}|\tb z_t,\tb w)p(\tb x|\tb z_1,\tb w).
            \end{align*}
            Then,
            \begin{align*}
                \mathcal{L}(\tb w)
                &= \mathbb{E}_{q(\tb z_{1:T}|\tb x)}\left[ \log{\frac{p(\tb z_T)\{\prod_{t=2}^{T}p(\tb z_{t-1}|\tb z_t,\tb w)\}p(\tb x|\tb z_1, \tb w)}{q(\tb z_1 |\tb x)\prod_{t=2}^{T}q(\tb z_t|\tb z_{t-1},\tb x)}} \right]
            \\  &= \mathbb{E}_{q(\tb z_{1:T}|\tb x)}\left[ \log{p(\tb z_T)} + \sum_{t=2}^{T}\log{\frac{p(\tb z_{t-1}|\tb z_t,\tb w)}{q(\tb z_t|\tb z_{t-1},\tb x)}} + \log{\frac{p(\tb x|\tb z_1,\tb w)}{q(\tb z_1|\tb x)}} \right].
            \end{align*}
            We use Bayes' Theorem and substitute $q(\tb z_t|\tb z_{t-1}, \tb x) = \frac{q(\tb z_{t-1}|\tb z_t, \tb x)q(\tb z_t|\tb x)}{q(\tb z_{t-1}|\tb x)}$.
            Then,
            \begin{align*}
                \mathcal{L}(\tb w)
                &= \mathbb{E}_{q(\tb z_{1:T}|\tb x)}\left[\sum_{t=2}^{T}\log{\frac{p(\tb z_{t-1}|\tb z_t,\tb w)}{q(\tb z_{t-1}|\tb z_t, \tb x)}} + \log{p(\tb x|\tb z_1, \tb w)} \right]
            \\  &= \int q(\tb z_1|\tb x)\log{p(\tb x|\tb z_1,\tb w)}d\tb z_1 - \sum_{t=2}^{T}\int \text{KL}(q(\tb z_{t-1}|\tb z_t,\tb x)\|p(\tb z_{t-1}|\tb z_t,\tb w))q(\tb z_t|\tb x)d\tb z_t
            \end{align*}
            The first term is called reconstruction term, and the second term is called consistency terms.

            We now substitute $q(\tb z_{t-1}|\tb z_t,\tb x)$ of consistency terms with conditional distribution in Theorem \ref{thm:revCondDist}.
            \[
            \text{KL}(q(\tb z_{t-1}|\tb z_t,\tb x)|p(\tb z_{t-1}|\tb z_t,\tb w)) = \frac{1}{2\beta_t}\|\tb m_t(\tb x,\tb z_t)-\bm\mu(\tb z_t,\tb w_t)\|^2 + \text{const}
            \]
            and approximate the reconstruction term with Monte Carlo method. Note that we do not need reparametrization method since the encoder is fixed.

            Now we can maximize the ELBO term about parameter $\tb w$.

        \subsubsection{Predicting the Noise instead of Mean}
            Modifying the ELBO so that we predict the added noise instead of the denoised data itself is known to produce higher quality results.
            We first rearrange the equation that can be derived from Theorem \ref{thm:diffKern} to give
            \[
            \tb x = \frac{1}{\sqrt{\alpha_t}}\tb z_t - \frac{\sqrt{1-\alpha_t}}{\sqrt{\alpha_t}}\bm\epsilon_t
            \]
            Using this, we can rewrite the mean $\tb m_t(\tb x,\tb z_t)$ as
            \[
            \tb m_t(\tb x,\tb z_t)=\frac{1}{\sqrt{1-\beta_t}}\left( \tb z_t = \frac{\beta_t}{\sqrt{1-\alpha_t}}\bm\epsilon_t \right)
            \]
            and similarly, we can rewrite $\bm\mu(\tb z_t, \tb w, t)$, introducing a neural network $\tb g(\tb z_t, \tb w, t)$ that predicts 
            total noise accumulation:
            \[
            \bm\mu(\tb z_t, \tb w, t)=\frac{1}{\sqrt{1-\beta_t}}\left( \tb z_t - \frac{\beta_t}{\sqrt{1-\alpha_t}}\tb g(\tb z_t, \tb w, t) \right)
            \]

            Substituting these eqautions to the consistency terms, we get
            \begin{align*}
                \text{KL}(q(\tb z_{t-1}|\tb z_t,\tb x)\|p(\tb z_{t-1}|\tb z_t,\tb w))
                &=\frac{\beta_t}{2(1-\alpha_t)(1-\beta_t)}\|\tb g(\tb z_t, \tb w, t)-\bm\epsilon_t\|^2 + \text{const}
            \\  &=\frac{\beta_t}{2(1-\alpha_t)(1-\beta_t)}\|\tb g(\sqrt{\alpha_t}\tb x + \sqrt{1-\alpha_t}\bm\epsilon_t, \tb w, t)-\bm\epsilon_t\|^2 + \text{const}
            \end{align*}

            The reconstruction term can be approximated with a sampled value of $\tb z_1$. Then we have
            \begin{align*}
                \log{p(\tb x|\tb z_1,\tb w)}
                &=-\frac{1}{2\beta_1}\| \tb x-\bm\mu(\tb z_1,\tb w,1) \|^2 + \text{const}
            \\  &=-\frac{1}{2(1-\beta_1)}\|\tb g(\tb z_1,\tb w,1)-\bm\epsilon_1\|^2 + \text{const}
            \end{align*}
            Using $\alpha_1=1-\beta_1$, we can combine the reconstruction term and consistency term:
            \[
            \mathcal{L}(\tb w)=-\sum_{t=1}^{T}\frac{\beta_t}{2(1-\alpha_t)(1-\beta_t)}\|\tb g(\sqrt{\alpha_t}\tb x + \sqrt{1-\alpha_t}\bm\epsilon_t,\tb w, t) - \bm\epsilon_t\|^2
            \]
            ignoring the constants.
            A further empirical study showed that performance is improved by ommiting the factor $\frac{\beta_t}{2(1-\alpha_t)(1-\beta_t)}$.

            \begin{algorithm}[H]
                \DontPrintSemicolon
                \SetKwInOut{Input}{Input}
                \SetKwInOut{Output}{Output}
                \SetArgSty{textnormal}

                \caption{Denoising diffusion probabilistic model (DDPM) training algorithm}

                \Input{training data $\data = \{\tb x_1,\dots,\tb x_N\}$ \newline
                noise schedule $\{\beta_1,\dots,\beta_T\}$}
                \Output{network parameters \tb w}

                \BlankLine
                \hrule
                
                \For{$t \in \{1,\dots,T\}$}
                {
                    $\alpha_t \leftarrow \prod_{\tau=1}^{t}(1-\beta_t)$\;
                }
                \Repeat{convergence}
                {
                    $\mathcal{L} \leftarrow 0$\;
                    \For{mini-batch $\{\tb x_n\}$}
                    {
                        $t \sim \{1,\dots,T\}$\;
                        $\bm\epsilon \sim \mathcal{N}(\bm 0, \tb I)$\;
                        $\tb z_t \leftarrow \sqrt{\alpha_t}\tb x + \sqrt{1-\alpha_t}\bm\epsilon$\;
                        $\mathcal{L}(\tb w) \leftarrow \mathcal{L}(\tb w) + \|\tb g(\tb z_t,\tb w, t)-\bm\epsilon\|^2$\;
                    }
                    $\tb w \leftarrow \tb w - \nabla_w\mathcal{L}(\tb w)$\;
                }

                \Return{\tb w}
            \end{algorithm}
    \subsection{Sampling}
        \begin{algorithm}[H]
            \DontPrintSemicolon
            \SetKwInOut{Input}{Input}
            \SetKwInOut{Output}{Output}
            \SetArgSty{textnormal}

            \caption{Denoising diffusion probabilistic model (DDPM) sampling algorithm}

            \Input{trained denoising network $\tb g(\tb z,\tb w,t)$ \newline
            noise schedule $\{\beta_1,\dots,\beta_T\}$}
            \Output{sample vector \tb x in data space}

            \BlankLine
            \hrule
            $\tb z_T \leftarrow \mathcal{N}(\tb 0, \tb I)$\;
            \For{$t \in \{T,\dots,2\}$}
            {
                $\alpha_t \leftarrow \prod_{\tau=1}^{t}(1-\beta_\tau)$\;
                $\bm\mu(\tb z_t, \tb w, t) \leftarrow \frac{1}{\sqrt{1-\beta_t}}(\tb z_t - \frac{\beta_t}{1-\alpha_t}\tb g(\tb z_t,\tb w,t))$\;
                $\bm\epsilon \leftarrow \mathcal{N}(\tb 0, \tb I)$\;
                $\tb z_{t-1} \leftarrow \bm\mu(\tb z_t,\tb w,t) + \sqrt{\beta_t}\bm\epsilon$\;
            }
            $\tb x \leftarrow \frac{1}{\sqrt{1-\beta_1}}(\tb z_1 - \frac{\beta_1}{\sqrt{1-\alpha_1}}\tb g(\tb z_1,\tb w,1))$\;
            \Return{\tb x}
        \end{algorithm}
\section{Score Matching}
\end{document}